\input{sol_pre.tex}
\input{q_2025-05-24_1.tex}
\begin{document}
\input{sol_head.tex}

\Savol{1}
EKUK eng kiçik bölişi uçun har bir tub köpaytuvçining darajasi $a$ va $b$ orasida imkon qadar teng bölinadi: $3=1+2$, $5=2+3$, $2=1+1$.
\[\text{EKUK}_{\min}=2^2\cdot3^3\cdot7=756.\]
\Javob{D}

\Savol{2}
$n-1$ soni $2\cdot3\cdot5\cdot7=210$ ga bölinadi: $n=210k+1$. \ $[200;1000]$ da: $211,\ 421,\ 631,\ 841$ — $4$ ta.
\Javob{B}

\Savol{3}
Har biri $n$ ta detal: \ $\dfrac n{20}-\dfrac n{21}=\dfrac16$ soat $\Rightarrow$ $\dfrac n{420}=\dfrac16$ $\Rightarrow$ $n=70$.
\Javob{B}

\Savol{4}
Asl bahosi $x$ va $y$ (mln): \ $1{,}4(x+y)=35$ $\Rightarrow$ $x+y=25$; \ $1{,}25x+1{,}5y=35$ $\Rightarrow$ $0{,}25y=35-31{,}25$ $\Rightarrow$ $y=15$.
\Javob{C}

\Savol{5}
$a=2^{49}$: \ $a^2+1=(a+1)^2-2a=(a+1)^2-\left(2^{25}\right)^2=\left(a+1-2^{25}\right)\left(a+1+2^{25}\right)$. Natija: $2^{49}-2^{25}+1$.
\Javob{B}

\Savol{6}
$7\mp\sqrt{24}=\left(\sqrt6\mp1\right)^2$: \ $\dfrac1{\sqrt6-1}-\dfrac1{\sqrt6+1}=\dfrac25$.
\Javob{A}

\Savol{7}
$2664=3\cdot888$, $1776=2\cdot888$: \ $\sqrt{888^2(9-4-1)}=2\cdot888=1776$.
\Javob{B}

\Savol{8}
$a_{98}=a_{99}$ $\Rightarrow$ $d=0$, barça hadlar $a$ ga teng: \ $a+1=a-2a+a=0$ $\Rightarrow$ $a=-1$. \ Yiğindi $100\cdot(-1)=-100$.
\Javob{C}

\Savol{9}
Progressiya $b,\ bq,\ bq^2,\ bq^3$: \ $b(1+q)=4$, \ $bq^2(1+q)=36$ $\Rightarrow$ $q=3$, $b=1$.\\
$C_1=x_1x_2=3$, \ $C_2=x_3x_4=9\cdot27=243$: \ $\sqrt{729}=27$.
\Javob{C}

\Savol{10}
$n=2024$: \ $(n+1)^3-(n-1)^3=6n^2+2=a$ $\Rightarrow$ $\sqrt{\dfrac{a-2}6}=n=2024$.
\Javob{B}

\Savol{11}
$\dfrac1{\sqrt x}-\dfrac1{\sqrt y}=\dfrac{\sqrt y-\sqrt x}{\sqrt{xy}}$, şuning uçun
\[\frac{\sqrt{xy}}{y\sqrt y-x\sqrt x}\cdot\frac{\sqrt y-\sqrt x}{\sqrt{xy}}=\frac{\sqrt y-\sqrt x}{\left(\sqrt y-\sqrt x\right)\left(y+\sqrt{xy}+x\right)}=\frac1{x+\sqrt{xy}+y}.\]
\Javob{A}

\Savol{12}
$\cos1125^\circ=\cos45^\circ=\frac{\sqrt2}2$, \ $\sin\frac{17\pi}6=\sin\frac{5\pi}6=\frac12$: \ $\dfrac{\sqrt2+1}2$.
\Javob{A}

\Savol{13}
$6\sin^2x-4\sin x\cos x=5\sin^2x+5\cos^2x$ $\Rightarrow$ $\tg^2x-4\tg x-5=0$ $\Rightarrow$ $\tg x=-1$ yoki $\tg x=5$.\\
$\tg x=-1$: $x=-\frac\pi4,\ -\frac{5\pi}4$; \ $\tg x=5$: $x=\arctg5-\pi$ ($\arctg5>\frac\pi4$, $\arctg5-2\pi<-\frac{3\pi}2$). Jami $3$ ta.
\Javob{B}

\Savol{14}
$\left(x^2-1\right)3^x-9\left(x^2-1\right)>0$ $\Rightarrow$ $\left(x^2-1\right)\left(3^x-9\right)>0$. Nollar $-1;\ 1;\ 2$:
\[x\in(-1;\,1)\cup(2;\,\infty).\]
\Javob{A}

\Savol{15}
Aniqlaniş sohasi: $x>2$, $x\ne3$, $(x-1)(x-3)>0$ $\Rightarrow$ $x>3$. U holda asos $x-2>1$:
\[x^2-4x+3<1\ \Rightarrow\ x^2-4x+2<0\ \Rightarrow\ x<2+\sqrt2.\quad x\in\left(3;\,2+\sqrt2\right).\]
\Javob{B}

\Savol{16}
$x^2-3=5$ $\Rightarrow$ $x=\pm2\sqrt2$; \ $x^2-3=-5$ — yeçim yöq.
\Javob{A}

\Savol{17}
$S=x+y+z$. Qöşamiz: $2S\cdot S=288$ $\Rightarrow$ $S=\pm12$. \ $x+y=\dfrac{108}S$, \ $z=S-\dfrac{108}S$:\\
$S=12$: $\dfrac9{3}=3$; \ $S=-12$: $\dfrac{-9}{-3}=3$.
\Javob{C}

\Savol{18}
Aniqlaniş sohasi: $x^2+3x-10\le0$ $\Rightarrow$ $x\in[-5;\,2]$. Ikkala köpaytuvçi manfiy emas, demak köpaytma $=0$:\\
$x=-5$, $x=2$ yoki $x^2-5x+2=0$ (butun ildizi yöq). \ Yiğindi $-5+2=-3$.
\Javob{C}

\Savol{19}
$2^{2x}-\frac32\cdot2^{2x}+2=2-\frac12\cdot4^x>0$ $\Leftrightarrow$ $x<1$. Kasr musbat: $x<1$ va $x>-4$ (ikkinçi holat $x>1$, $x<-4$ bölmaydi).\\
$x\in(-4;\,1)$: $-3-2-1+0=-6$.
\Javob{C}

\Savol{20}
$y-3=\dfrac x{1+x}$ $\Rightarrow$ $x=\dfrac{y-3}{4-y}$. \ $f^{-1}(x)=\dfrac{x-3}{4-x}=\dfrac{3-x}{x-4}$.
\Javob{D}

\Savol{21}
$f(-x)=-f(x)$: \ $f(x)\left(1+x^2\right)=x^3+5x$ $\Rightarrow$ $f(1)=\frac62=3$, \ $f(-1)=-3$.
\Javob{B}

\Savol{22}
$t=\sqrt{2x-4}$, $x=\frac{t^2+4}2$: \ $x\pm2t=\dfrac{(t\pm2)^2}2$. \ $x>4$ da $t>2$: \ $f(x)=\dfrac{(t+2)-(t-2)}{\sqrt2}=2\sqrt2$ — özgarmas.\\
Demak $f'(6)=0$.
\Javob{B}

\Savol{23}
\[\int_0^{\pi/4}\tg^6x\cdot\frac{dx}{\cos^2x}=\int_0^{\pi/4}\tg^6x\,d(\tg x)=\frac{\tg^7x}7\Big|_0^{\pi/4}=\frac17.\]
\Javob{B}

\Savol{24}
Kvadrat tomoni $2r$: \ $\dfrac{8r}{2\pi r}=\dfrac4\pi$.
\Javob{A}

\Savol{25}
$\angle AOB=90^\circ+\frac{\angle C}2=150^\circ$ $\Rightarrow$ $\angle C=120^\circ$. \ $S=\frac12\cdot13\cdot8\cdot\sin120^\circ=26\sqrt3$.
\Javob{A}

\Savol{26}
$\triangle ABD\sim\triangle CBA$ ($\angle B$ umumiy, $\angle BAD=\angle C$), öxşaşlik koeffitsiyenti $\dfrac{AD}{CA}=\dfrac23$.\\
$S_{ABD}=\frac49S_{ABC}$ $\Rightarrow$ $S_{ADC}=\frac59S_{ABC}$. \ Nisbat $\dfrac45$.
\Javob{B}

\Savol{27}
$DG=s$. Uçburçakni tomonlari $16+\frac s2$ va $16$ bölgan töğri törtburçak içiga olamiz va burçakdagi uçta töğri burçakli uçburçakni ayiramiz ($w=16+\frac s2$):
\[S=16w-\frac12\cdot8\cdot16-\frac12(w-8)(16-s)-\frac12ws=8w-4s=128.\]
(Natija $s$ ga boğliq emas.)
\Javob{B}

\Savol{28}
Tomon $1$. $HC\parallel AB$, $\angle KCB=45^\circ$, $\angle CKB=\angle ABK=67{,}5^\circ$ $\Rightarrow$ $CK=CB=1$, \ $HC=1+\sqrt2$, \ $HK=\sqrt2$.\\
$S_1=\frac12\cdot1\cdot\frac{\sqrt2}2=\frac{\sqrt2}4$. \ $F$ dan $HC$ gaça masofa $1+\frac{\sqrt2}2$:
\[S_2=S_{HKF}+S_{HFG}=\frac12\sqrt2\left(1+\frac{\sqrt2}2\right)+\frac12\cdot1\cdot\frac{\sqrt2}2=\frac{2+3\sqrt2}4,\qquad\frac{S_2}{S_1}=\frac{2+3\sqrt2}{\sqrt2}=3+\sqrt2.\]
\Javob{D}

\Savol{29}
$S_{ABCD}=\frac{15}2$ (koordinatalar böyiça). $AE$ yuzni teng bölsa, $S_{AED}=\frac{15}4=\frac12\cdot4\cdot y_E$ $\Rightarrow$ $y_E=\frac{15}8$.\\
$DC$ töğri çiziği: $x=4-\frac y3$ $\Rightarrow$ $x_E=4-\frac58=\frac{27}8$.
\Javob{B}

\Savol{30}
Asos tomoni $a$: $h=\frac a2$, apofema $\frac a{\sqrt2}$. \ $S_{\text{yon}}=2a\cdot\frac a{\sqrt2}=\sqrt2a^2=144\sqrt2$ $\Rightarrow$ $a=12$, $h=6$.\\
$V=\frac13\cdot144\cdot6=288$.
\Javob{C}

\Savol{31}
$D\cup E=\{3;4;6;7\}$. \ $(A\cup B)\cap C=\{2;3;4;6;7;8\}\cap\{3;\dots;7\}=\{3;4;6;7\}$.
\Javob{C}

\Savol{32}
$\{0;1;2;3;4\}$ dan $3$ ta raqam tanlanadi, tartib yagona: $C_5^3=10$.
\Javob{C}

\newpage
\Savol{33}
Kub qirrasi $2$. \ $OM=ON=MN=\sqrt2$ ($O$ dan qirra örtasigaça — yoq diagonalining yarmi), \ $MD_1=\sqrt5$, \ $OD_1=\sqrt3$, \ $CM=CN=\sqrt5$.
\fig{\begin{tikzpicture}[scale=1.5,font=\small]
\coordinate (A) at (0,0); \coordinate (B) at (2,0); \coordinate (C) at (2.8,0.6); \coordinate (D) at (0.8,0.6);
\coordinate (A1) at (0,2); \coordinate (B1) at (2,2); \coordinate (C1) at (2.8,2.6); \coordinate (D1) at (0.8,2.6);
\coordinate (M) at (0.4,0.3); \coordinate (N) at (1,0); \coordinate (O) at (1.4,1.3);
\draw (A)--(B)--(C)--(C1)--(D1)--(A1)--cycle; \draw (A1)--(B1)--(C1); \draw (B)--(B1);
\draw[dashed] (A)--(D)--(C); \draw[dashed] (D)--(D1);
\draw[axs] (M)--(O)--(N)--cycle; \draw[ax] (M)--(D1)--(O); \draw[ax] (M)--(C)--(N);
\fill (M) circle (1.3pt); \fill (N) circle (1.3pt); \fill (O) circle (1.3pt);
\node[below left] at (A) {$A$}; \node[below right] at (B) {$B$}; \node[right] at (C) {$C$}; \node[above left] at (D) {$D$};
\node[left] at (A1) {$A_1$}; \node[above left] at (B1) {$B_1$}; \node[above right] at (C1) {$C_1$}; \node[above left] at (D1) {$D_1$};
\node[left] at (M) {$M$}; \node[below] at (N) {$N$}; \node[right] at (O) {$O$};
\end{tikzpicture}}

\textbf{33.} $\triangle MON$ teng tomonli: $\angle MON=60^\circ$, \ $\sin60^\circ=\dfrac{\sqrt3}2$ — \textbf{D}.

\textbf{34.} $\cos\angle MD_1O=\dfrac{5+3-2}{2\sqrt5\cdot\sqrt3}=\dfrac3{\sqrt{15}}=\dfrac{\sqrt{15}}5$ — \textbf{A}.

\textbf{35.} $\cos\angle MCN=\dfrac{5+5-2}{2\cdot5}=\dfrac45$ — \textbf{E}.
\Javob{33 — D, \ 34 — A, \ 35 — E}

\newpage
\Savol{36}
$y=\dfrac{x^3+3}4$ belgilaymiz. U holda tenglama $y^3=4x-3$ ga, belgilaş esa $x^3=4y-3$ ga teng. Ayiramiz:
\[x^3-y^3=4(y-x)\ \Rightarrow\ (x-y)\left(x^2+xy+y^2+4\right)=0\ \Rightarrow\ x=y.\]
($x^2+xy+y^2+4>0$.) Demak $x^3-4x+3=0$ $\Rightarrow$ $(x-1)\left(x^2+x-3\right)=0$:
\[x_1=1,\qquad x_{2,3}=\frac{-1\pm\sqrt{13}}2.\]
a) $3$ ta ildiz. \quad b) $1+(-1)=0$.
\Javob{a) 3; \quad b) 0}

\newpage
\Savol{37}
$x$ ga nisbatan kvadrat tenglama: $D=16\cos^2(xy)-16\ge0$ $\Rightarrow$ $\cos^2(xy)=1$.

$\cos(xy)=1$: \ $(x+2)^2=0$, $x=-2$; \ $xy=2\pi k\in[-\pi;\pi]$ $\Rightarrow$ $xy=0$, $y=0$.

$\cos(xy)=-1$: \ $(x-2)^2=0$, $x=2$; \ $xy=\pm\pi$ $\Rightarrow$ $y=\pm\dfrac\pi2$.

Juftliklar: $(-2;\,0)$, $\left(2;\,\frac\pi2\right)$, $\left(2;\,-\frac\pi2\right)$.

a) $3$ ta. \quad b) $x+y$: $-2$; \ $2+\frac\pi2$; \ $2-\frac\pi2>0$. Eng kiçigi $-2$.
\Javob{a) 3; \quad b) $-2$}

\newpage
\Savol{38}
$x=-1$: \ $2f(-1)+3f(-1)+f(-1)=-1$ $\Rightarrow$ $f(-1)=-\frac16$, ya'ni $3f(-1)=-\frac12$:
\[2f(x)+f\left(\tfrac1x\right)=x+\tfrac12.\]
a) $x=1$: \ $3f(1)=\frac32$ $\Rightarrow$ $f(1)=\dfrac12$.

b) $x$ ni $\frac1x$ ga almaştiramiz: $2f\left(\frac1x\right)+f(x)=\frac1x+\frac12$. Sistemadan:
\[3f(x)=2x+1-\frac1x-\frac12\ \Rightarrow\ f(x)=\frac13\left(2x-\frac1x+\frac12\right).\]
$f(x)=1$: \ $2x-\frac1x-\frac52=0$. \ $u=\frac1x$ uçun: $2u^2+5u-4=0$, \ $u_1+u_2=-\frac52$, \ $u_1u_2=-2$:
\[\frac1{x_1^2}+\frac1{x_2^2}=u_1^2+u_2^2=\frac{25}4+4=\frac{41}4.\]
\Javob{a) $\frac12$; \quad b) $\frac{41}4$}

\newpage
\Savol{39}
Urinma $(1;\,2)$ va $(3;\,0)$ nuqtalardan ötadi: $k=\dfrac{0-2}{3-1}=-1$. \ Demak $f(1)=2$, $f'(1)=k=-1$.

a) $f'(1)=-1$.

b) $g(x)=x^3-2x^2+x+3$: \ $g(1)=3$, \ $g'(x)=3x^2-4x+1$, \ $g'(1)=0$.
\[h'(1)=g'(1)f(1)+g(1)f'(1)=0\cdot2+3\cdot(-1)=-3.\]
\Javob{a) $-1$; \quad b) $-3$}

\newpage
\Savol{40}
Grafikdan: $f(-4)=4$, $f(-1)=-2$, $f(3)=6$; $f$ \ $(-4;-1)$ da kamayadi, $(-1;3)$ da ösadi.

a) $f'\ge0$ bölgan joyda integral ostidagi ifoda $0$, $f'<0$ da esa $-2f'$:
\[\int_{-4}^{-1}(-2f'(x))\,dx=-2\bigl(f(-1)-f(-4)\bigr)=-2(-2-4)=12.\]
b) $f+xf'=(xf)'$, \ $ff'=\left(\frac{f^2}2\right)'$:
\[\left[xf(x)+\frac{f^2(x)}2\right]_{-4}^{3}=\left(18+18\right)-\left(-16+8\right)=44.\]
\Javob{a) 12; \quad b) 44}

\newpage
\Savol{41}
\fig{\begin{tikzpicture}[scale=1.5,font=\small]
\coordinate (B) at (0,1.591); \coordinate (D) at (1.555,0.334); \coordinate (A) at (-1.0,-1.237); \coordinate (C) at (3.599,1.591); \coordinate (P) at (0.778,0.962);
\draw (0,0) circle (1.591); \draw (A)--(B)--(C)--cycle; \draw (B)--(D);
\draw[ax] (A)--(P);
\node[above] at (B) {$B$}; \node[below right] at (D) {$D$}; \node[below left] at (A) {$A$}; \node[above right] at (C) {$C$}; \node[aux,above right] at (P) {$P$};
\end{tikzpicture}}
Urinma va vatar orasidagi burçak: $\angle CBD=\angle BAD$. $\angle C$ umumiy $\Rightarrow$ $\triangle CBD\sim\triangle CAB$, koeffitsiyent $\dfrac{BD}{AB}=\dfrac23$.

a) $S_{CBD}=\frac49S_{CAB}$ $\Rightarrow$ $S_{ABD}=\frac59S_{CAB}$, \ $S_{CBD}=\frac45S_{ABD}$.\\
$ABD$ teng yonli: $AP$ — balandlik ($P$ — $BD$ örtasi), $AP=\sqrt{9-1}=2\sqrt2$, \ $S_{ABD}=\frac12\cdot2\cdot2\sqrt2=2\sqrt2$.
\[S_{BDC}=\frac45\cdot2\sqrt2=\frac{8\sqrt2}5.\]
b) Aylana $ABD$ uçburçakka taşqi çizilgan:
\[R=\frac{AB\cdot BD\cdot AD}{4S_{ABD}}=\frac{3\cdot2\cdot3}{8\sqrt2}=\frac{9\sqrt2}8.\]
\Javob{a) $\frac{8\sqrt2}5$; \quad b) $\frac{9\sqrt2}8$}

\newpage
\Savol{42}
\fig{\begin{tikzpicture}[scale=0.8,font=\small]
\coordinate (A) at (0,0); \coordinate (B) at (5.565,0); \coordinate (C) at (5.565,5.565); \coordinate (D) at (0,5.565); \coordinate (E) at (2.783,2.064); \coordinate (P) at (2.783,0); \coordinate (Q) at (5.565,2.064);
\fill[okgreen!30] (A)--(E)--(C)--cycle;
\draw (A)--(B)--(C)--(D)--cycle; \draw (A)--(C); \draw (A)--(E)--(C); \draw (E)--(B);
\draw[ax] (E)--(P); \draw[ax] (E)--(Q);
\node[below left] at (A) {$A$}; \node[below right] at (B) {$B$}; \node[above right] at (C) {$C$}; \node[above left] at (D) {$D$}; \node[above left] at (E) {$E$}; \node[aux,below] at (P) {$P$}; \node[aux,right] at (Q) {$Q$};
\end{tikzpicture}}
$AE=BE$ $\Rightarrow$ $E$ — $AB$ ning örta perpendikulyarida. Tomon $a$, $EP=t$ ($P$ — $AB$ örtasi), $EQ\perp BC$: $EQ=\frac a2$, $QC=a-t$.
\[t^2=12-\frac{a^2}4,\qquad 20=\frac{a^2}4+(a-t)^2=a^2-2at+12\ \Rightarrow\ t=\frac{a^2-8}{2a}.\]
a) $\dfrac{\left(a^2-8\right)^2}{4a^2}=\dfrac{48-a^2}4$ $\Rightarrow$ $\left(a^2\right)^2-32a^2+32=0$ $\Rightarrow$ $a^2=16\pm4\sqrt{14}$.\\
$E$ uçburçak içida: $t>0$ $\Rightarrow$ $a^2>8$. \ $S=a^2=16+4\sqrt{14}=4\left(4+\sqrt{14}\right)$.

b) $S_{ACE}=S_{ABC}-S_{ABE}-S_{BCE}=\dfrac{a^2}2-\dfrac{at}2-\dfrac{a^2}4=\dfrac{a^2}2-\dfrac{a^2-8}4-\dfrac{a^2}4=2$.
\Javob{a) $4\left(4+\sqrt{14}\right)$; \quad b) 2}

\newpage
\Savol{43}
\fig{\begin{tikzpicture}[scale=1.15,font=\small]
\coordinate (A) at (1.6,0); \coordinate (B) at (1.131,1.131); \coordinate (C) at (0,1.6); \coordinate (D) at (-1.131,1.131); \coordinate (E) at (-1.6,0); \coordinate (F) at (-1.131,-1.131); \coordinate (G) at (0,-1.6); \coordinate (H) at (1.131,-1.131); \coordinate (K) at (-2.623,1.749);
\draw (A)--(B)--(C)--(D)--(E)--(F)--(G)--(H)--cycle; \draw (A)--(K)--(H);
\draw[ax] (A)--(H); \draw[ax] (0,0) circle (1.6);
\draw[axs] ($(A)!0.22cm!(K)$) -- ($(A)!0.22cm!(K)!0.22cm!90:(A)$) -- ($(A)!0.22cm!(H)$);
\node[right] at (A) {$A$}; \node[above right] at (B) {$B$}; \node[above] at (C) {$C$}; \node[above] at (D) {$D$}; \node[left] at (E) {$E$}; \node[below left] at (F) {$F$}; \node[below] at (G) {$G$}; \node[below right] at (H) {$H$}; \node[above left] at (K) {$K$};
\end{tikzpicture}}
Sakkizburçakka taşqi aylana çizamiz. $\angle HAD$ — içki çizilgan, $HGFED$ yoyiga ($4\cdot45^\circ=180^\circ$) tiraladi: $\angle HAK=90^\circ$.\\
Içki burçak $\angle AHG=135^\circ$ $\Rightarrow$ $\angle AHK=135^\circ-60^\circ=75^\circ$.

a) $AH=HK\cos75^\circ=12\cdot\dfrac{\sqrt6-\sqrt2}4=3\left(\sqrt6-\sqrt2\right)=6\sqrt{2-\sqrt3}$.

b) $S_{AKH}=\frac12\cdot HK\cos75^\circ\cdot HK\sin75^\circ=\frac{144}4\sin150^\circ=18$.
\Javob{a) $6\sqrt{2-\sqrt3}$; \quad b) 18}

\newpage
\Savol{44}
Asos tomonlari: $\frac{\sqrt3}4a^2=9\sqrt3$ $\Rightarrow$ $a=6$; \ $b=12$. Içki çizilgan aylana radiuslari $r_1=\frac a{2\sqrt3}=\sqrt3$, \ $r_2=2\sqrt3$.

Öq va yon yoq apofemasi orqali kesim — töğri burçakli trapetsiya $O_1M_1M_2O_2$ ($O_1M_1=r_1$, $O_2M_2=r_2$); şar katta doirasi unga içki çizilgan.
\fig{\begin{tikzpicture}[scale=0.9,font=\small]
\coordinate (O2) at (0,0); \coordinate (M2) at (3.464,0); \coordinate (O1) at (0,4.899); \coordinate (M1) at (1.732,4.899); \coordinate (Q) at (0,2.449); \coordinate (T) at (1.732,0);
\draw (O2)--(M2)--(M1)--(O1); \draw[ax] (O1)--(O2); \draw[okgreen!60!black,thick] (Q) circle (2.449); \draw[ax] (M1)--(T);
\fill (Q) circle (1.5pt);
\node[below] at (O2) {$O_2$}; \node[below right] at (M2) {$M_2$}; \node[above] at (O1) {$O_1$}; \node[above right] at (M1) {$M_1$}; \node[left] at (Q) {$Q$}; \node[aux,below] at (T) {$T$};
\end{tikzpicture}}
Bir nuqtadan urinmalar tengligidan: apofema $l=M_1M_2=r_1+r_2=3\sqrt3$. \ $h^2=l^2-(r_2-r_1)^2=27-3=24$, \ $R=\frac h2=\sqrt6$.

a) $S_{\text{yon}}=3\cdot\dfrac{6+12}2\cdot3\sqrt3=81\sqrt3$.

b) $S_{\text{şar}}=4\pi R^2=4\cdot3\cdot6=72$.
\Javob{a) $81\sqrt3$; \quad b) 72}

\newpage
\Savol{45}
a) $EAB$ va $EBC$ yoqlarni bir tekislikka yoyamiz: $\angle AEC=60^\circ+60^\circ=120^\circ$, $EF=2$, $EC=4$.
\fig{\begin{tikzpicture}[scale=0.8,font=\small]
\coordinate (E) at (0,0); \coordinate (A) at (-4,0); \coordinate (B) at (-2,-3.464); \coordinate (C) at (2,-3.464); \coordinate (F) at (-2,0);
\draw (E)--(A)--(B)--(C)--cycle; \draw (E)--(B); \draw[axs] (F)--(C); \fill (F) circle (1.5pt);
\node[above] at (E) {$E$}; \node[above left] at (A) {$A$}; \node[below] at (B) {$B$}; \node[below right] at (C) {$C$}; \node[above] at (F) {$F$};
\end{tikzpicture}}
\[FC^2=2^2+4^2-2\cdot2\cdot4\cos120^\circ=28\ \Rightarrow\ FC=2\sqrt7.\]
(Asos orqali yöl: $FC^2=3^2+\left(4+\sqrt3\right)^2=28+8\sqrt3>28$ — uzunroq.)

b) $t=\dfrac{2\sqrt7}{\sqrt7}\cdot2=4$ daqiqa.
\Javob{a) $2\sqrt7$; \quad b) 4}

\ShowJavoblar
\end{document}
